% Source: https://artofproblemsolving.com/community/c3h1605901_random_algebra_question
% Stable identity: aops-topic-1605901-q01
\textbf{The other root.} Since the coefficients are rational and $\sqrt7$ is irrational, the conjugate $\dfrac{1-\sqrt7}{2}$ is also a root. (Indeed, if $r=\dfrac{1+\sqrt7}{2}$ is a root, then dividing $2022x^2+mx+n$ by the rational quadratic $(x-r)\left(x-\bar r\right)=x^2-x-\dfrac{3}{2}$ leaves a remainder $px+q$ with rational $p,q$ that vanishes at the irrational $r$, which forces $p=q=0$.)

\textbf{Vieta's formulas.} The roots have sum $1$ and product $\dfrac{1-7}{4}=-\dfrac{3}{2}$. Hence
\[
-\dfrac{m}{2022}=1,\qquad\dfrac{n}{2022}=-\dfrac{3}{2},\qquad\text{so}\qquad m=-2022,\quad n=-3033 ,
\]
and $m-n=-2022+3033=\boxed{1011}$.
